\chapter{Fresh reconstruction of the physical model and conditional assembly}
\label{component:p9b00698ff36a}
\noindent\textbf{Component }\path{p9b00698ff36a}\par
{\small\noindent Source: \path{REC/routing/fresh_physical_normalization_conditional_assembly.tex}\par}
{\small\noindent SHA256: \path{9b00698ff36aab2f233a0d2932a439ad7fb145384d10545302ecb1441513394d}\par}
{\small\noindent Inventory use: Fresh physical real-qubit/finite-convexification model, exact normalization and separated matrix-pairing assembly | RFA physical-reduction/normalization/closure portions including RFA10--RFA12\par}
\medskip
\section{Fresh reconstruction of the physical model and conditional assembly}
\label{p9b00698ff36a:sec:fresh-r21-conditional-assembly}

\paragraph{Status of this source.}
This is a newly derived proof, not a transcription of the missing final
R21 manuscript. The physical-model theorem, the homogeneous normalization
identity, the separated conic/convex closure lemma, and the matrix-pairing
lemma below are proved here. The final existence implication is expressly
conditional on the local-certificate hypothesis (LC). Neither a historical
admission nor a remembered checksum is used to discharge (LC).

\subsection{An exact finite-dimensional physical model}

Let $Q\subset\mathbb R^8$ be the set of correlations obtained from arbitrary
finite-dimensional complex bipartite tensor-product states and two binary
measurements per party. The coordinates are
\[
 (a_0,a_1,b_0,b_1,E_{00},E_{01},E_{10},E_{11}).
 \tag{RFA1}
\]
A binary measurement is represented by an effect $0\preceq E\preceq I$,
with observable $2E-I$. Put
$Z=\operatorname{diag}(1,-1)$ and
$X=\left(\begin{smallmatrix}0&1\\1&0\end{smallmatrix}\right)$.
For $(c_A,s_A),(c_B,s_B)\in S^1$ and $\psi\in S^3\subset\mathbb R^4$,
let $\Phi$ be the eight expectations of
\[
 A_0=B_0=Z,\qquad A_1=c_AZ+s_AX,\qquad B_1=c_BZ+s_BX
 \tag{RFA2}
\]
in the state $\psi$. Every coordinate of $\Phi$ is polynomial in these
eight real parameters.

\begin{theorem}[Fresh physical reduction]
With $\mathcal Z=S^1\times S^1\times S^3$,
\[
 Q=\operatorname{conv}\Phi(\mathcal Z)
   =\left\{\sum_{i=1}^{9}w_i\Phi(z_i):
       z_i\in\mathcal Z,\ w_i\geq0,\ \sum_iw_i=1\right\}.
 \tag{RFA3}
\]
In particular $Q$ is compact, convex and semialgebraic over the real field.
\end{theorem}

\begin{proof}
First, both effects of a party can be dilated with the same embedding.
For an effect $E$ on $\mathcal H$ define, on $\mathcal H\oplus\mathcal H$,
\[
 \widehat P_E=
 \begin{pmatrix}E&\sqrt{E(I-E)}\\
                 \sqrt{E(I-E)}&I-E\end{pmatrix},
 \qquad Vh=(h,0).
 \tag{RFA4}
\]
The square root commutes with $E$ by its spectral construction. Direct
block multiplication gives $\widehat P_E^2=\widehat P_E$ and
$V^*\widehat P_EV=E$. Thus for both measurement settings the same $V$
compresses the two projections to the original effects. Applying this
construction to each party and replacing $\rho$ by
$(V_A\otimes V_B)\rho(V_A\otimes V_B)^*$ preserves every marginal and
joint expectation. We may therefore assume the two observables of each
party are Hermitian involutions.

Here is the needed two-dimensional block decomposition. If $U,V$ are
Hermitian involutions, $C=UV$ is unitary and $UCU=C^*$. For an eigenvector
$e$ of $C$ with nonreal eigenvalue $\lambda$, $Ue$ has eigenvalue
$\overline\lambda$. The vectors are orthogonal and their span is invariant
under $U$ and $V=UC$. Its orthogonal complement is invariant too. Repeating
on orthogonal eigenvectors gives two-dimensional blocks. On the real
eigenvalue spaces $\lambda=\pm1$ one has $V=\pm U$; diagonalizing $U$ gives
one-dimensional blocks. This proves the decomposition in finite dimension
without an assumption on multiplicities.

On a two-dimensional block $U$ has eigenvalues $1,-1$. In its eigenbasis,
the other nontrivial involution has the form
$\left(\begin{smallmatrix}c&\overline z\\z&-c\end{smallmatrix}\right)$,
where $c$ is real and $c^2+|z|^2=1$. A diagonal change of phase leaves
$U=Z$ and makes $z$ real. Thus the pair has the form $Z,cZ+sX$.
A one-dimensional block with outcomes $a_0,a_1\in\{\pm1\}$ embeds into
this form: use the $a_0$ eigenvector of $Z$ and take
$c=a_0a_1,s=0$. The unused basis vector does not affect the experiment.

Apply both local decompositions. All measurement operators are block
diagonal, so their expectations depend only on the diagonal blocks of
the state. The traces of these positive blocks are nonnegative weights
summing to one. Every nonzero block, after normalization and the harmless
one-dimensional embedding, is a two-qubit state with real measurement
matrices. Decompose that state into finitely many complex pure states.
For a real symmetric matrix $M$ and a vector $r+is$,
\[
 (r+is)^*M(r+is)=r^TMr+s^TMs,\qquad
 \|r+is\|^2=\|r\|^2+\|s\|^2.
 \tag{RFA5}
\]
The cross terms cancel by symmetry. The nonzero real and imaginary parts,
normalized separately, are real pure states with their squared norms as
weights. This proves $Q\subseteq\operatorname{conv}\Phi(\mathcal Z)$.

Conversely each atom is a physical two-qubit experiment. A finite convex
combination is realized by a finite direct sum on each party and a state
with the corresponding classical diagonal weights in matching summands.
Therefore the reverse inclusion holds.

Any convex combination of more than nine points in $\mathbb R^8$ can be
shortened. The augmented vectors $(1,q_i)\in\mathbb R^9$ are linearly
dependent, so there are $\alpha_i$, not all zero, with
$\sum_i\alpha_i=0$ and $\sum_i\alpha_iq_i=0$. Change their common sign if
needed so that some $\alpha_i>0$, and replace the positive weights $w_i$
by $w_i-\tau\alpha_i$, where
$\tau=\min_{\alpha_i>0}w_i/\alpha_i$. They remain nonnegative, keep the
same sum and point, and at least one vanishes. Iterate. Zero weights
allow a fixed list of nine atoms. The right side of (RFA3) is the
continuous image of the compact product of a simplex and nine copies of
$\mathcal Z$, hence is compact. It is semialgebraic by the projection
theorem for real closed fields applied to its displayed finite polynomial
formula. Convexity has already been proved.
\end{proof}

For any real closed extension $K\supset\mathbb R$, define $Q_K$ by
(RFA3) with all variables in $K$. Define
$P=\{f\in\mathbb R[q]_{\leq1}:f\geq0\text{ on }Q\}$.
The defining formula
$\forall q\,(q\in Q\Rightarrow f(q)\geq0)$, together with quantifier
elimination and elementary transfer, gives exactly
\[
 P_K=\{f\in K[q]_{\leq1}:f\geq0\text{ on }Q_K\}.
 \tag{RFA6}
\]
Compactness was used only over $\mathbb R$. No compactness assertion for
the order topology of a non-Archimedean field is needed.

\subsection{Separated fields and the full homogeneous parent}

Fix arbitrary real closed extensions $K_c,K_t\supset\mathbb R$ and put
$T=K_c\otimes_{\mathbb R}K_t$. The factors are labelled; they may be two
copies of the same extension. There is no assumption of algebraic
disjointness and no chosen ordering on $T$. For
$f=f_0+\sum_{\nu=1}^8f_\nu q_\nu$ define
\[
 \mathscr S(f,q)=f_0\otimes1+\sum_{\nu=1}^8f_\nu\otimes q_\nu.
 \tag{RFA7}
\]

The event and CHSH convention in the recovered physical core is
\[
 S=E_{00}-E_{01}+E_{10}+E_{11},\quad H=1-S/2,\quad
 f=dH+4\sum_{\xi}\xi\,p_\xi,
 \quad \xi\in\{\lambda,n,m,C,A,D,\alpha,\beta\}.
 \tag{RFA8}
\]
Its original full parent is recovered as actual source text, not merely
a remembered identity. We give its homogeneous form explicitly.

Choose projective angle pairs $(a_0,a_1)$ and $(b_0,b_1)$ in $K_t^2$,
each nonzero, and set
$\Delta_A=a_0^2+a_1^2>0$, $\Delta_B=b_0^2+b_1^2>0$.
The corresponding circle coordinates are
\[
 c_A=(a_0^2-a_1^2)/\Delta_A,\quad s_A=2a_0a_1/\Delta_A,
 \tag{RFA9}
\]
and similarly for Bob. These cover every circle point over $K_t$:
if $c\ne-1$ use $(1+c,s)$, and if $c=-1$ use $(0,1)$.
No affine half-angle chart is therefore needed at infinity.
For $\Psi=(W,X,Y,Z)$ put
\[
 \begin{split}
 L_m&=a_0b_1W-a_0b_0X+a_1b_1Y-a_1b_0Z,\\
 L_D&=a_1b_0W+a_1b_1X-a_0b_0Y-a_0b_1Z.
 \end{split}
\]
Products of a coefficient-field element with a testing-field expression
in the next identity mean their separated tensor product. The complete
homogeneous parent is
\[
 \begin{split}
 \mathcal F={}&4\lambda\Delta_A(b_0W+b_1X)^2
       +4n\Delta_A\Delta_BX^2+4mL_m^2\\
 &+4C\Delta_B(a_1W-a_0Y)^2+4A\Delta_A\Delta_BY^2\\
 &+(4D+2d)L_D^2
       +(4\alpha+2d)\Delta_A(b_1Y-b_0Z)^2\\
 &+(4\beta+2d)\Delta_B(a_0X+a_1Z)^2
       -2d\Delta_A\Delta_BZ^2.
 \end{split}                                                    \tag{RFA10}
\]
For a nonzero state let $N=\|\Psi\|^2>0$ and let $z$ be the atom
with the circle coordinates (RFA9) and unit state $\Psi/\sqrt N$.
Then
\[
 \mathcal F=(1\otimes\Delta_A\Delta_BN)\,
                 \mathscr S(f,\Phi(z)).
 \tag{RFA11}
\]

To check the normalization, first take $a_0=b_0=1$, $u=a_1$ and
$v=b_1$. The plus/minus measurement eigenvectors are
$(1,u)/\sqrt{1+u^2}$ and $(u,-1)/\sqrt{1+u^2}$.
Their product amplitudes give, in the order of (RFA8),
\[
 \frac{(W+vX)^2}{1+v^2},\quad X^2,\quad
 \frac{(vW-X+uvY-uZ)^2}{(1+u^2)(1+v^2)},\quad
 \frac{(uW-Y)^2}{1+u^2},\quad Y^2,
\]
\[
 \frac{(uW+uvX-Y-vZ)^2}{(1+u^2)(1+v^2)},\quad
 \frac{(vY-Z)^2}{1+v^2},\quad
 \frac{(X+uZ)^2}{1+u^2}.
 \tag{RFA12}
\]
For a normalized state these are exactly the eight selected event
probabilities. For a general state they are multiplied by $N$.
Substitution of the matrices (RFA2) into $H$ gives its numerator
\[
 2\{(uW+uvX-Y-vZ)^2+(1+u^2)(vY-Z)^2
       +(1+v^2)(X+uZ)^2-(1+u^2)(1+v^2)Z^2\}.
\]
This proves (RFA11) in the affine chart. Homogenization gives (RFA10);
it is a polynomial identity and therefore also gives the zero and
infinite projective coordinates without dividing by $a_0$ or $b_0$.
The zero state has zero parent separately.

If $\mathcal F$ or $(s\otimes1)\mathcal F$ has an exact finite SOS
certificate, with $s\in K_c^{>0}$, the original slack is restored by
multiplying each root by
\[
 s^{-1/2}\otimes(\Delta_A\Delta_BN)^{-1/2}.
 \tag{RFA13}
\]
All these inverses are taken within one field before forming a tensor.
The identity does not permit cancellation of an arbitrary mixed tensor.
Additional positive source gauges are handled in the same way. A state
chart or a change of frame must first prove its own denominator nonzero.

\paragraph{Literal commuting tests.}
These require no limiting local certificate. If one party's two
involutions commute, on a simultaneous eigenspace the CHSH operator
$S=A_0(B_0-B_1)+A_1(B_0+B_1)$ is either $\pm2B_0$ or $\pm2B_1$.
Thus $H=I-S/2$ is positive semidefinite in the testing field. For
$d\geq0$, its contribution is a pure-source square times a
pure-testing square. Every nonnegative event packet is also separated
SOS. This covers the actual commuting angle axes and infinities for
every nonnegative canonical coefficient tuple. It does not cover merely
infinitesimal nonzero angles, which remain in (LC).

\subsection{The closed faces and the surviving source reduction}

Two recovered sources supply actual mathematical inputs:
\path{physical_core.tex} gives the canonical event basis, the
determinant/coercivity argument, the finite-contact extreme-source
reduction and the nine-term source decomposition;
\path{closed_faces_short.tex} gives the two explicit matrix lifts
for the omitted exclusive pairs. Their exact recovered paths and hashes
are recorded in the reconstruction map. Their roles are not supplied
by the historical R21 review narrative.

We state the precise source-generation conclusion used below. Every
transferred valid affine functional is either an event sum or a finite
nonnegative conic combination of the following types, with finite real
Bell relabelings restored:
\begin{itemize}
\item event sources and noncoercive closed-face sources;
\item physical C1 sources with $\lambda=n=0$, $d,m,C,A,\alpha>0$,
      $D,\beta\geq0$, $D+\beta>0$, and an actual finite positive
      noncommuting stationary contact;
\item physical C2 sources with $\lambda=\beta=0$,
      $d,m,C,A,D,\alpha>0$, $n\geq0$, and the same actual-contact
      condition.
\end{itemize}
The source sum is formed before choosing $d=1$ or $C=2$.
This is a source-field assertion; source contacts are not testing points.

For completeness, the logical content of the recovered reduction is as
follows. A functional nonnegative on the nonsignaling polytope is a
nonnegative combination of its sixteen event inequalities. Otherwise it
is negative at a PR vertex, because deterministic vertices are physical.
The explicit nonsignaling vertex classification and the nine evaluations
in the recovered event table give (RFA8) with nonnegative event
coefficients and $d>0$. Normalize their total sum to one. The physical
coefficient section is compact and convex in a space of affine dimension
at most eight and hence is a convex combination of at most nine extreme
points. The recovered determinant has nonnegative outer Laurent
coefficients. Its coercive cases have a finite positive minimum; the
printed opposite-perturbation argument makes an extreme minimum zero.
The four kernel and two angular equations leave a nonradial perturbation
if more than six event coefficients are active. The complementary state
pivots, exact adjugate identity and convex Laurent profile preserve
physicality under that perturbation. Coercivity requires at least five
active events by the explicit exclusive-pair classification. Compatible
omissions have precisely the two C1/C2 orientations above; every
noncoercive support lies on one of the two closed faces. The finite
source decomposition and these contact assertions have fixed first-order
formulas, so their real proofs transfer. This is the actual content
required of the recovered source proof; merely naming a C1 or C2 row
would not suffice.

Here is a direct separated-tensor consequence of the recovered
closed-face lifts. Write a given lift, including all free nonnegative
source-coordinate diagonal blocks, as the homogeneous matrix
\[
 M(c,y)=\sum_\nu c_\nu A_\nu+\sum_jy_jB_j\succeq0.
\]
The supplied matrices have an explicit strict feasible point in the
free source coordinates and auxiliaries. If a linear functional
$\ell(c)$ is nonnegative on the projected cone, there exists $Y\succeq0$
with
\[
 \operatorname{tr}(A_\nu Y)=\ell_\nu,\qquad
 \operatorname{tr}(B_jY)=0.                          \tag{RFA14}
\]
One can prove this without leaving an attainment issue implicit.
For the linear map $M$, its adjoint image $M^*(\mathrm{PSD})$ is
closed: at a fixed strict point $z_*$, $M(z_*)\succeq\epsilon I$,
so convergence of $M^*(Y_i)$ for $Y_i\succeq0$ bounds
$\epsilon\operatorname{tr}(Y_i)$ and thus bounds $Y_i$.
A convergent subsequence supplies the limiting witness. Separation
of this closed cone proves that its dual is exactly
$\{z:M(z)\succeq0\}$, yielding (RFA14).

Apply this over the real field with $\ell$ equal to a testing
evaluation. The fixed existential statement for $Y$ transfers to $K_t$;
source feasibility transfers independently to $K_c$. For a source
witness $X=M(c,y)\succeq0$, factor $X=BB^T$ in $K_c$ and
$Y=CC^T$ in $K_t$. Positive semidefinite matrices over a real closed
field have such finite factorizations, including singular matrices,
by symmetric diagonalization or the pivoted quadratic-form algorithm.
Then
\[
 \sum_{i,k}X_{ik}\otimes Y_{ik}
  =\sum_{a,b}\left(\sum_iB_{ia}\otimes C_{ib}\right)^2.
 \tag{RFA15}
\]
The adjoint equations (RFA14) cancel the auxiliary variables and identify
the left side with the original slack. This proves the closed-face
tensor conclusion from the actual explicit face lifts.

Event slacks also need no local certificate: for a real unit atom,
the joint event projector $R$ satisfies
$p=\psi^TR\psi=\|R\psi\|^2$ over $K_t$. Its entries are testing-field
elements, so this is a finite pure-testing SOS. A nonnegative source
weight contributes its own square root. Constants and the zero
functional are included because outcomes of one input pair sum to one.

\subsection{The missing local premise, stated without a status substitution}

\paragraph{Local-certificate hypothesis (LC).}
For every pair of real closed extensions $K_c,K_t$, every physical
canonical C1 or C2 source in the source-generation statement above,
and every independent projective angle pair and testing state, the
complete original parent (RFA10) has a finite SOS certificate in $T$,
possibly after multiplication by a proved strictly positive pure-source
factor. Every source coefficient face and gap, all testing signs,
all state coordinates, and all projective endpoints must be included.
When local charts are used, their union and restoration must be proved.

The reconstruction does not yet establish (LC). In particular the
missing positive-scale completion, original weak-scale entry/scoping
proofs, uniform-C1 local replacements and the C14 and weak-C2 local
closures cannot be replaced by their historical admission summaries.
The accompanying map separates exact surviving inputs, alternate
candidate proofs, fragmentary statements and absent certificate bodies.

One sufficient way to prove (LC) is to restore a finite implication
tree, valid for every actual source and test, whose leaves are
proved full-parent certificates or proved impossible domains.
The branch predicates must carry the actual source equations and
literal boundaries. This description is a proof obligation, not an
assertion that the missing tree has been recovered.

\subsection{Conditional conic/convex closure and finite-lift implication}

\begin{lemma}[Fresh separated assembly]
Let $f=\sum_{j=1}^r a_jg_j$ in the source field, with $a_j\geq0$,
and let $q=\sum_{i=1}^s w_iq_i$ in the testing field, with
$w_i\geq0$ and $\sum_iw_i=1$. If every elementary slack has a
finite certificate
$\mathscr S(g_j,q_i)=\sum_kh_{ijk}^2$ in $T$, then
\[
 \mathscr S(f,q)=
 \sum_{i,j,k}\bigl((\sqrt{a_j}\otimes\sqrt{w_i})h_{ijk}\bigr)^2.
 \tag{RFA16}
\]
\end{lemma}
\begin{proof}
Affineness in the testing point, including its constant term, gives
$\mathscr S(f,q)=\sum_i(1\otimes w_i)\mathscr S(f,q_i)$.
Linearity in the source gives
$\mathscr S(f,q_i)=\sum_j(a_j\otimes1)\mathscr S(g_j,q_i)$.
Substitute the finite certificates. The required nonnegative square
roots exist in the separate real closed fields. Zero weights give
zero roots; no weight is inverted. This proves the displayed identity
in the tensor ring, without identifying its two factors.
\end{proof}

\begin{theorem}[Fresh conditional universal theorem]
Assume (LC), with the actual source reduction and explicit closed-face
lifts specified above. Then for all real closed $K_c,K_t\supset\mathbb R$,
all $f\in P_{K_c}$ and all $q\in Q_{K_t}$,
$\mathscr S(f,q)$ is a finite sum of squares in
$K_c\otimes_{\mathbb R}K_t$. Consequently $Q$ has a finite
semidefinite lift.
\end{theorem}
\begin{proof}
Use the finite source decomposition. Event and closed-face slacks
were proved above. For every other source atom apply (LC), then
(RFA11)--(RFA13) to restore its exact affine slack.
Use the nine-atom testing decomposition (RFA3) and the assembly
lemma. Real/complex states, mixed states, the zero functional and
zero weights were already included.

For the final implication take two labelled copies of an arbitrary
single real closed extension $K$. The set and valid-affine cone are
exactly their transfers by (RFA3) and (RFA6). The resulting SOS is
in $K\otimes_{\mathbb R}K$. The factor flip converts the
coefficient-first convention to the point-first convention of
Claus Scheiderer's Corollary 3.19, which applies to a closed convex
semialgebraic set in arbitrary finite ambient dimension. It proves
finite semidefinite-lift existence.
\end{proof}

\paragraph{Exact boundary of the new conclusion.}
The last theorem is conditional until actual proofs establish (LC).
No unconditional finite-lift conclusion is claimed by this reconstruction.
Even after (LC) is proved, this criterion gives existence, not evaluated
matrices, a numerical size bound, or a common separation-rank bound.
The explicit feature, affine-pencil and coefficient-validation items
E1--E3 remain subsequent constructive tasks.

\paragraph{Primary external criterion.}
C. Scheiderer, \emph{Second-order cone representation for convex sets
in the plane}, arXiv:2004.04196v2, Corollary 3.19 and its proof,
printed page 15, \url{https://arxiv.org/pdf/2004.04196}.
The primary text has been reopened during this reconstruction.
Theorem 3.10 in the same source distinguishes a common tensor-rank
bound from unrestricted finite SOS. Neither local square counts nor
the number of routing branches supplies that stronger bound.
